\section{Introduction}
\label{sec:introduction}

The Arnold fixed-point problem asks how much of the critical-point theory of
a small autonomous Hamiltonian persists for an arbitrary Hamiltonian path.
One proposed lower bound is the stable Morse number: the least number of
critical points after allowing auxiliary quadratic variables. We construct a
family for which this bound fails by a fixed proportion, while the ambient
manifolds remain simply connected and Kähler of real dimension twenty-two.

Let $W$ be a connected closed smooth manifold. Its \emph{stable Morse number}
$\SM(W)$ is the minimum of $\#\Crit(F)$ over integers $k\geq0$ and smooth
Morse functions $F:W\times\R^k\to\R$ such that, for $k>0$,
\[
 F(x,z)=Q(z)\qquad\text{for all sufficiently large }|z|,
\]
where $Q$ is a nondegenerate quadratic form of arbitrary signature. The
radius is uniform in $x$. For $k=0$ there is no condition at infinity; the
minimum over this subclass is the ordinary Morse number $\Morse(W)$.
In particular, $\SM(W)\leq\Morse(W)$.

For a connected closed symplectic manifold $(M,\omega)$ and
$H\in C^\infty((\R/\Z)\times M,\R)$, write $\phi_H^t$ for the Hamiltonian
flow, with convention $\omega(X_{H_t},\cdot)=dH_t$. A fixed point of
$\phi_H^1$ is \emph{nondegenerate} if its derivative has no eigenvalue $1$.
We write
\[
 \Fix_0(\phi_H^1;H)
 =\{x\in\Fix(\phi_H^1):t\mapsto\phi_H^t(x)
          \text{ is contractible in }M\}.
\]
This is a set of initial points; choices of capping disks do not change its
cardinality. The unrestricted stable Morse formulation asks whether
$\#\Fix_0(\phi_H^1;H)\geq\SM(M)$ whenever all fixed points are
nondegenerate.

\begin{theorem}\label{thm:main}
For every integer $m\geq1$, there exist a simply connected closed Kähler
manifold $(M_m,\omega_m)$ of real dimension $22$ and a smooth one-periodic
Hamiltonian $H_m$ such that all fixed points of $\phi_{H_m}^1$ are
nondegenerate and
\begin{align*}
 \SM(M_m)&=80+1968m,\\
 \#\Fix(\phi_{H_m}^1)
 &=\#\Fix_0(\phi_{H_m}^1;H_m)=80+1952m.
\end{align*}
Consequently, with $\delta=1/128$, for every $R>0$ some member of this
family satisfies
\[
 \SM(M_m)\geq R,\qquad
 \#\Fix(\phi_{H_m}^1)\leq(1-\delta)\SM(M_m).
\]
\end{theorem}

\subsection*{Context and earlier work}
Arnold's questions connect Hamiltonian dynamics to the topology of smooth
functions~\cite{Arnold1986}. Conley and Zehnder established the torus case by
variational methods~\cite{ConleyZehnder1983}. Floer's construction of a
chain complex from pseudoholomorphic curves established a new approach to
homological fixed-point bounds~\cite{Floer1989}. Later extensions include the
weakly monotone theory of Hofer--Salamon and Ono and the virtual-moduli
methods of Fukaya--Ono, Liu--Tian, and Ruan
\cite{HoferSalamon1995,Ono1995,FukayaOno1999,LiuTian1998,Ruan1999}.
The resulting rational homological lower bound and the geometric Morse
bounds are different assertions: integral torsion can force more critical
points than the total rational Betti number.

More recent work gives lower bounds over arbitrary fields and over the
integers~\cite{AbouzaidBlumberg2021,BaiXu2022,BaiXu2026}. The integral
Floer bound uses homology with a cyclic grading determined by the first
Chern class. The stable Morse number retains the integer grading, and
these yield different bounds for the examples below; the comparison is
computed in Remark~\ref{counts:integral-Floer-bound}.

The stable Morse formulation is stated explicitly in
\cite[Conjecture~1.2]{Golovko2020}. Damian showed that ordinary and stable
Morse numbers need not agree~\cite{Damian2002}. Positive results for the
stable bound include the theorem of Dimitroglou Rizell and Golovko for
generic Hamiltonians when
both the symplectic area class and the first Chern class vanish on
$\pi_2(M)$~\cite[Corollary~1.4]{DRG2017}; P\"oder Balkest\aa hl obtained an
independent proof using the simple homotopy type of the Hamiltonian Floer
complex~\cite[Corollary~1.2]{PoderBalkestahl2021}. These hypotheses do not
hold for the projective blow-ups used here.
The almost-quadratic convention used in these comparisons gives the same
stable Morse number as our exact condition at infinity; see
Remark~\ref{morse:infinity-conventions}.

The companion manuscript~\cite{OpenAICompanion} gives a simply connected
Kähler counterexample using mixed-prime torsion and a Hamiltonian
involution. The present construction develops that mechanism in a fixed
dimension and produces an arbitrarily large stable Morse number with a
uniform relative deficit. We reproduce the geometric and dynamical
arguments needed for the construction. The principal extra step is to
embed arbitrarily many disjoint copies of suitable centers in $\CP^9$:
the number of centers may grow without increasing their dimensions or the
dimension of the ambient projective space.

\subsection*{Mechanism and proof structure}
All homology groups below have integral coefficients unless another ring
is displayed. For a finite abelian group $T$, let $d(T)$ denote its minimum
number of generators. Put
\begin{equation}\label{intro:lambda}
\begin{gathered}
 b(W)=\sum_i\rank H_i(W;\Z),\qquad
 t_i(W)=d(\Tor H_i(W;\Z)),\\
 \lambda(W)=b(W)+2\sum_i t_i(W).
\end{gathered}
\end{equation}
The geometric Morse inequalities give $\SM(W)\geq\lambda(W)$ even with
auxiliary quadratic variables. For simply connected manifolds in the
dimensions used here, Smale's handle theorem supplies an ordinary Morse
function attaining this bound~\cite[Theorem~2.4 and the following
paragraph]{Smale1963}. Section~\ref{sec:morse} proves the stable lower
bound and states the equality; Appendix~\ref{app:handles} recalls its
handle-theoretic proof.

The generator count combines different primes by a maximum rather than
a sum. For example,
\[
 d\bigl((\Z/2)^a\oplus(\Z/3)^q\bigr)=\max(a,q).
\]
We will arrange that the prime contributing the larger multiplicity
changes with the homological degree. Redistributing those degrees by
taking products can then change the generator count in a way that the
free ranks do not detect.

To arrange the needed homology, take two projective surfaces whose torsion
groups are respectively $(\Z/2)^2$ and $\Z/3$ in degrees $1$ and $2$.
The blow-up centers are the first surface and the product of the second
with $\CP^1\times\CP^1$, of complex dimensions $2$ and $4$.
Section~\ref{sec:geometry} embeds $m$ copies of each disjointly into
$\CP^9$ and blows them up to obtain a simply connected Kähler manifold
$X_m$. The projection argument applies more generally to every finite
collection of smooth projective varieties of dimension at most $d$, using
one ambient $\CP^{2d+1}$.

The dynamical construction uses the Kähler surface $Y$ obtained by blowing
up the four coordinate corners of $\CP^1\times\CP^1$. The half-turn of the
diagonal circle action fixes exactly the four exceptional projective lines.
Thus its product with the identity on $X_m$ has four fixed components on
$M_m=X_m\times Y$, each diffeomorphic to $C_m=X_m\times\CP^1$.
Section~\ref{sec:dynamics} perturbs this involution so that each component
contributes precisely the critical points of a chosen Morse function on
$C_m$. A short-period argument excludes every other fixed point. The
construction therefore realizes $4\lambda(C_m)$ nondegenerate fixed points.

Products with $\CP^1$ and $Y$ combine the torsion in nearby degrees with
different multiplicities. Taking the maximum of the two primary
multiplicities in each degree then gives different savings in the number
of generators. Section~\ref{sec:counts} computes
\[
 \lambda(X_m\times Y)-4\lambda(X_m\times\CP^1)=16m
\]
and completes the proof. This separates the stable Morse count from the
Hamiltonian count without any distinction between ordinary and stable
Morse theory on the simply connected manifolds in the construction.
